% Independent mathematical derivation of incoming 02_RESIDUE_BOUNDARY_VALUE.md.
% Inclusion fragment. The fixed-minor receiving map is proved for every chart.
\begingroup
\section{An explicit reflected trace and both original boundary metrics}
\label{sec:RBV}

Keep the original fixed hypothetical simple quartet, its full arithmetic
unit, actual period, logarithm branch and conductor. Write
\[
\begin{gathered}
a\equiv1\pmod4,\quad a\ge17,\quad c=a/2,\quad
q=(a+1)^2,\quad q'=(a-7)^2,\quad m=16a-48-v,\quad0\le v\le80,\\
\lambda_{u,b}=c+(2u-a)\delta+i(2b-a)\gamma,\qquad
0<\delta<\tfrac12,\quad\gamma>2,\quad0\le u,b\le a,\\
\chi(S)=\prod_{u,b=0}^a(S-\lambda_{u,b}),\quad
E=\mathbb C[S]/(\chi),\quad
\mathcal P_{q-1}\xrightarrow{\ \sim\ }E .
\end{gathered}\tag{RBV1}
\]
The last map sends an actual polynomial to its residue class, with
the unique degree-below-$q$ remainder as inverse. Roots are distinct
by equality of real and imaginary parts. No packet factors or coordinate
phases are discarded.

\subsection{The explicit polynomial and its exact trace}
Put
\[
t_0=\frac\pi{2\gamma},\qquad r_0=t_0\delta,\qquad
c_0=8\log4-10>0,\qquad
P_a(S)=\sum_{j=0}^{4a}\frac{t_0^{2j}(S-c)^{2j}}{(2j)!}.
\tag{RBV2}
\]
Its degree is $8a\le m-1$ since $8a\ge129\ge49+v$.
Also $8a<q$. Define the original finite reflected form
\[
\mathcal W(f,g)=\sum_{\lambda}\overline{f(a-\overline\lambda)}g(\lambda).
\tag{RBV3}
\]
The reflection permutes the complete grid and is an involution.
Changing its index in the conjugate of $\mathcal W(g,f)$ proves
that this form is Hermitian. The polynomial $P_a$ has real even
coefficients in $S-c$, so
$\overline{P_a(a-\overline\lambda)}=P_a(\lambda)$.
Hence $\mathcal W(P_a,P_a)=\sum_\lambda P_a(\lambda)^2$, a real
number because conjugation also permutes the complete grid.

For $z=\lambda_{u,b}-c$, the actual comparison value is
\[
\cosh(t_0z)=i\sin\!\left(\frac{\pi(2b-a)}2\right)
                   \sinh(r_0(2u-a)).\tag{RBV4}
\]
The sine is $+1$ or $-1$, and $2u-a$ is a nonzero odd integer.
Moreover
$t_0|z|\le at_0\sqrt{\delta^2+\gamma^2}<2a$:
$\delta/\gamma<1/4$ and $\pi<22/7$ give the final strict bound.
For $x\ge0$ and an integer $n\ge1$, the exponential series gives
$\sum_{j\ge n}x^j/j!\le x^ne^x/n!$, because
$(n+j)!\ge n!j!$. The elementary integral comparison
$\log n!=\sum_{j=1}^n\log j\ge\int_1^n\log x\,dx
=n\log n-n+1$ gives $n!\ge(n/e)^n$. Therefore
\[
\begin{split}
|P_a(c+z)-\cosh(t_0z)|
&\le\frac{e^{2a}(2a)^{8a+1}}{(8a+1)!}\\
&\le e^{2a}(e/4)^{8a+1}\le e^{-c_0a}=:\epsilon_a.
\end{split}\tag{RBV5}
\]
This bounds the full omitted even tail by a larger exponential tail.
Impose the explicit eventual guard $\epsilon_a\le r_0/4$.
Since $|\sinh(r_0(2u-a))|\ge r_0$, write the actual value as
$is+e$ with $s$ real and $|e|\le|s|/4$. Direct expansion gives
\[
\Re(is+e)^2\le-s^2+2|s||e|+|e|^2\le-7s^2/16.
\tag{RBV6}
\]
The elementary finite geometric sum gives exactly
\[
\begin{gathered}
\sum_{u=0}^a\cosh(2r_0(2u-a))
=\frac{\sinh(2r_0(a+1))}{\sinh(2r_0)},\\
\sum_{\lambda}\cosh(t_0(\lambda-c))^2
=-\frac{a+1}{2}
 \left[\frac{\sinh(2r_0(a+1))}{\sinh(2r_0)}-(a+1)\right].
\end{gathered}\tag{RBV7}
\]
Consequently the incoming finite sign and its constants are correct:
\[
\boxed{\mathcal W(P_a,P_a)\le-N_a<0,\qquad
N_a=\frac{7(a+1)}{32}
\left[\frac{\sinh(2r_0(a+1))}{\sinh(2r_0)}-(a+1)\right].}
\tag{RBV8}
\]
The bracket is $2\sum_u\sinh^2(r_0(2u-a))>0$.
There is a stronger explicit finite lower bound. For real $s$ and
$0\le\epsilon\le|s|/2$, maximizing $\Re(is+x+iy)^2$ over
$x^2+y^2\le\epsilon^2$ gives $-(|s|-\epsilon)^2$.
For $s>0$, substitute $x^2\le\epsilon^2-y^2$: the remaining
concave quadratic $-s^2-2sy+\epsilon^2-2y^2$ takes its maximum
on $[-\epsilon,\epsilon]$ at $y=-\epsilon$ because its vertex
is $-s/2\le-\epsilon$. The case $s<0$ follows by reflection.
Consequently RBV8 also holds with
\[
\begin{split}
N_a^{\rm sharp}
&=(a+1)\sum_{u=0}^a
 (|\sinh(r_0(2u-a))|-\epsilon_a)^2\\
&=(a+1)\left\{\frac12
 \left[\frac{\sinh(2r_0(a+1))}{\sinh(2r_0)}-(a+1)\right]
-2\epsilon_a\frac{\cosh(r_0(a+1))-1}{\sinh r_0}
 +(a+1)\epsilon_a^2\right\}.
\end{split}\tag{RBV8a}
\]
The additional finite sum follows by summing the positive and
negative odd-index hyperbolic sines separately. In particular the
coefficient $7/32$ may be strengthened to $9/32$ on the same guard.
The original smaller $N_a$ is retained in the downstream displayed
bounds, all of which also hold with $N_a^{\rm sharp}$.
The exact error from RBV7 is at most
\[
2\epsilon_a(a+1)\sum_{u=0}^a|\sinh(r_0(2u-a))|
                         +q\epsilon_a^2.\tag{RBV9}
\]
Since the sum is bounded by $(a+1)e^{r_0a}$, this error divided by
the magnitude in RBV7 is
$O_{\delta,\gamma}(a e^{-(c_0+r_0)a})$.
Keeping also the lower exponential and $(a+1)$ term in RBV7 proves
\[
\mathcal W(P_a,P_a)
=-\frac{(a+1)e^{2r_0(a+1)}}{4\sinh(2r_0)}[1+o(1)].
\tag{RBV10}
\]

\subsection{Original arithmetic density and finite full-form constants}
The source remains
\[
\begin{gathered}
v_h(s)=2\xi(s)/h(s),\quad
w_h(y)=|v_h(1/2+iy)|^2/(2\pi),\quad
d\mu_a(y)=w_h^{*a}(y)\,dy,\\
M_h(\theta)=\int_{\mathbb R}e^{\theta y}w_h(y)\,dy,
\quad\alpha=\pi/2,\quad
d\sigma(y)=|\Gamma(1/4+iy/2)|^2dy/(2\pi),\quad
\sigma(\mathbb R)=\sqrt{2\pi}.
\end{gathered}\tag{RBV11}
\]
RTM1--3/GTM1 proves evenness, positivity almost everywhere and the
finite transforms at $|\theta|\le\alpha$, using the complete unit.
For the simple quartet, the original global constants in GTM1 are
\[
\begin{gathered}
B=50,\quad p=7/2,\quad
\vartheta_h=\int_{-1}^1w_h(y)\,dy,\quad\mathfrak M=M_h(\alpha),\\
b_a=\frac{c_h\vartheta_h^{a-3}e^{-\alpha(a-3)}(a-2)^{-50}}
 {C_\Gamma2^{25}},\qquad
A_a=\frac{C_h^{\rm tilt}}{c_\Gamma}a^{p+1}\mathfrak M^{a-1},\\
b_a(1+y^2)^{-25}d\sigma\le d\mu_a\le A_a d\sigma,\qquad
c_\Gamma=\sqrt2e^{-7/6},\quad C_\Gamma=\sqrt{2\sqrt5}e^{2/3}.
\end{gathered}\tag{RBV12}
\]
The constants $c_h,C_h^{\rm tilt}$ are exactly the original
TW7/BT bounds retained in RTM2, including both half-lines. This proof
does not posit a positive lower bound for the one-factor zeta numerator.

Here is the finite polynomial comparison needed for the whole
boundary bands. In the original orthonormal Gamma frame,
multiplication by $y$ has adjacent entries
$\sqrt{(n+1)(n+1/2)}\le n+1$ (GTM7).
On polynomials of degree at most $D$, its two shifts each have norm
at most $D+1$. Thus
$\|yF\|_\sigma\le2(D+1)\|F\|_\sigma$.
For $F\ne0$, divide the positive integral $|F|^2d\sigma$ by its
own total only to apply the following scalar inequality; the source
measure itself is never changed. Convexity of $(1+x)^{-25}$ on
$x\ge0$, proved by its positive second derivative, gives its tangent
lower bound at the mean of $y^2$. Integration therefore gives
\[
\|F\|_{\mu_a}^2\ge
b_a[1+4(D+1)^2]^{-25}\|F\|_\sigma^2.
\tag{RBV13}
\]
For all degrees through $2q$ we may use the exact constants
\[
L_a^-=b_a[1+4(2q+1)^2]^{-25},\qquad L_a^+=A_a,
\qquad\log(L_a^+/L_a^-)=O_h(a+\log q).\tag{RBV14}
\]
The original Gamma recurrence and mass dictionary are the
Meixner--Pollaczek formulas in DLMF (18.19.8)--(18.19.9), (18.22.8)
\cite{human:dlmf18}, with the exact substitution and mass verified
in LT2 and NTG3. This is a full coefficient-form comparison, so it
applies before either original boundary minimization.

\subsection{Concentration of the complete relation bands}
Set $Y=2^{-16}q$ and $R=a\sqrt{\delta^2+\gamma^2}$, and retain the
explicit eventual guard $R\le Y/4$.
For every polynomial $P$ of degree at most $q$ and
$F(S)=\chi(S)P(S)$, the original roots give
\[
|\chi(c+iy)|\le(Y+R)^q\quad(|y|\le Y),\qquad
|\chi(c+iy)|\ge(q-R)^q\quad(q\le y\le2q).
\tag{RBV15}
\]
Use the shifted Legendre basis
$\phi_d(y)=\sqrt{(2d+1)/q}\,L_d(2y/q-3)$ on $[q,2q]$,
$0\le d\le q$. Its Rodrigues formula and squared norm are the
Legendre entries of DLMF (18.5.5), Tables 18.5.1 and 18.3.1
\cite{human:dlmf18}. The full finite identity
\[
L_d(z)=2^{-d}\sum_{j=0}^d\binom dj^2(z-1)^{d-j}(z+1)^j
\]
gives $|L_d(z)|\le10^d$ on $|z|\le4$, using
$\sum_j\binom dj^2=\binom{2d}{d}\le4^d$.
For $|y|\le Y$, expansion of the original polynomial
$P(c+iy)$ in this basis and Cauchy--Schwarz prove
\[
|P(c+iy)|^2\le\frac{(q+1)^2}{q}100^q
                         \int_q^{2q}|P(c+it)|^2dt.
\tag{RBV16}
\]
Indeed the basis-square sum is bounded by
$q^{-1}100^q\sum_{d=0}^q(2d+1)=q^{-1}100^q(q+1)^2$.
The lower Gamma density on $[q,2q]$ is
$c_\Gamma e^{-\pi q}(1+2q)^{-1/2}$. Multiplying RBV15--16,
integrating the full central interval using mass $\sqrt{2\pi}$,
and using RBV14 in the denominator proves
\[
\begin{gathered}
\int_{|y|\le Y}|F(c+iy)|^2d\mu_a
 \le\bar\eta_a\|F\|_{\mu_a}^2,\qquad
\bar\eta_a=\min(1,\eta_a),\\
\eta_a=\frac{L_a^+}{L_a^-}\frac{\sqrt{2\pi}}{c_\Gamma}
 \frac{(q+1)^2\sqrt{1+2q}}q e^{\pi q}100^q
 \left(\frac{Y+R}{q-R}\right)^{2q},\\
\log\eta_a\le-c_1q+O_h(a+\log q),\qquad
c_1=30\log2-\pi-\log100>0.
\end{gathered}\tag{RBV17}
\]
The last ratio is at most $2^{-15}$ on the stated guard. This
bound applies to every vector of either complete relation band,
including arbitrary linear combinations and all original cross terms.

\subsection{Both boundary energies, including strict positivity}
Let $e_j=\chi\pi_j/\sqrt{\nu_j}$, where $\pi_j$ is the actual
monic polynomial orthogonal for the full weight
$|\chi(c+iy)|^2d\mu_a$ in the original coordinate. Define
\[
\mathcal B_0=\operatorname{span}(e_0,\ldots,e_{q-1}),\qquad
\mathcal B_1=\operatorname{span}(e_1,\ldots,e_q),\qquad
\Delta_iP=-\operatorname{Proj}_{\mathcal B_i}P.
\tag{RBV18}
\]
These are the original COB two windows. Each band consists of
$\chi$ times polynomials of degree at most $q$, so RBV17 applies
to every unit vector in it.
Their exact link to the original minimum forms follows without a
choice of diagonal trial metric. Let $s_N:E\to\mathcal P_N$ send
a class to its unique minimum-norm representative in $d\mu_a$.
For the original representative $P\in\mathcal P_{q-1}$, its entire
affine fibre is $P+\chi\mathcal P_{N-q}$. Completing its orthogonal
square gives
$s_NP=P-\operatorname{Proj}_{\chi\mathcal P_{N-q}}P$.
Consequently
$\Delta_0=s_{2q-1}-s_{q-1}$ and
$\Delta_1=s_{2q}-s_q$, exactly as in RBV18.
Writing $G_N=s_N^*s_N$ in the unchanged class frame, orthogonality
of these two nested projection differences gives
$G_{q-1}-G_{2q-1}=\Delta_0^*\Delta_0$ and
$G_q-G_{2q}=\Delta_1^*\Delta_1$.
Thus each coefficient matrix is exactly
$\Gamma_{i,I}=L_I^*\Delta_i^*\Delta_iL_I$ on its stated original
residue injection. This proves the map into both boundary energies,
including all cross terms of the actual minimum.
These are precisely COB.5--6 and COB.13 in
\path{COUPLED_ORIGINAL_CONE_BOUNDARY.tex}. The general positive
coefficient-square operation is the Schur minimum in Boyd and
Vandenberghe \cite[Appendix A.5.5, p.~650]{human:boyd-vandenberghe};
the preceding formulas prove its full original-fibre specialization.

In fact both $\Delta_i$ are injective on the entire original
$\mathcal P_{q-1}$. Pair the roots with indices $u$ and $a-u$.
Since $(a+1)^2/2$ is even, their original factors give exactly
\[
\chi(c+iy)=\prod_{b=0}^a\prod_{0\le u<a/2}
\left[(y-(2b-a)\gamma)^2+(2u-a)^2\delta^2\right]>0.
\tag{RBV19}
\]
If $P\perp\mathcal B_0$, test its orthogonality against
$\chi P\in\mathcal B_0$ to obtain
$\int\chi(c+iy)|P(c+iy)|^2d\mu_a=0$, hence $P=0$.
If $P\perp\mathcal B_1$, choose the unique $b$ such that
$F=P-b\chi$ is orthogonal to $e_0$. Since $\chi$ is a multiple
of $e_0$, $F$ remains orthogonal to $e_1,\ldots,e_q$.
Its degree is at most $q$, so $\chi F$ belongs to
$\operatorname{span}(e_0,\ldots,e_q)$. Testing against it and using
RBV19 gives $F=0$. Then $P=b\chi$ and $\deg P<q$ force $P=b=0$.
This proves strict positivity without an assumed missing metric bound.

On the original source line,
$|P_a(c+iy)|\le\cosh(t_0|y|)\le e^{t_0|y|}$.
The exact convolution transform, with all masses retained, is
\[
\int e^{\theta y}d\mu_a=M_h(\theta)^a.
\]
Evenness and $2t_0<\alpha$ therefore give, with
$\theta_1=(\alpha-2t_0)/2>0$, $\theta_2=2t_0+\theta_1<\alpha$,
\[
\|P_a\|_{\mu_a}^2\le2M_h(2t_0)^a,\qquad
\|1_{|y|>Y}P_a\|_{\mu_a}^2
 \le2e^{-\theta_1Y}M_h(\theta_2)^a.
\tag{RBV20}
\]
For a unit $f\in\mathcal B_i$, split its inner product with $P_a$
over the two original integration regions. Cauchy--Schwarz and RBV17
give
$|\langle f,P_a\rangle|\le\sqrt{\bar\eta_a}\|P_a\|+
\|1_{|y|>Y}P_a\|$. Taking the supremum over the full unit sphere
of the band, then $(x+y)^2\le2x^2+2y^2$, proves
\[
\boxed{0<\|\Delta_iP_a\|_{\mu_a}^2\le B_a,
\quad B_a=4\bar\eta_aM_h(2t_0)^a+
4e^{-\theta_1 2^{-16}q}M_h(\theta_2)^a,\quad i=0,1.}
\tag{RBV21}
\]
No band-dimension multiplier enters this estimate.

\subsection{Relative eigenvalue and evaluated discrepancy}
Let $G_i^E$ be the Gram of $\Delta_i$ on the actual monomial frame
of $\mathcal P_{q-1}$; it is strictly positive by RBV19.
Let $W^E$ be the matrix of RBV3 in that same frame. Then
\[
T_i^E=(G_i^E)^{-1/2}W^E(G_i^E)^{-1/2}
\]
is Hermitian. The coefficient column of $P_a$ gives the actual
Rayleigh quotient, so
\[
\lambda_{\min}(T_i^E)\le-N_a/B_a<0,\qquad
\|I-T_i^E\|\ge1+N_a/B_a.\tag{RBV22}
\]
This follows by diagonalizing the Hermitian matrix and expressing
the Rayleigh quotient as a positive weighted average of its real
eigenvalues. With the original two matrices, the discrepancy is
\[
\begin{split}
\|\Delta_iP_a\|_{\mu_a}^2-\mathcal W(P_a,P_a)
&=\frac{(a+1)e^{2r_0(a+1)}}{4\sinh(2r_0)}[1+o_h(1)],\\
\lim_{a\to\infty}\frac{\|\Delta_iP_a\|^2-\mathcal W(P_a,P_a)}
 {a e^{\pi\delta a/\gamma}}
&=\frac{e^{\pi\delta/\gamma}}{4\sinh(\pi\delta/\gamma)}.
\end{split}\tag{RBV23}
\]
Indeed RBV21 has logarithm at most
$-c_*q+O_h(a+\log q)$, where
\[
c_*=\min\{30\log2-\pi-\log100,
                   2^{-16}(\pi/2-\pi/\gamma)/2\}>0.
\tag{RBV24}
\]
This is negligible compared with RBV10. Also
$\log N_a=2r_0a+\log a+O_h(1)$, so RBV22 proves
\[
\liminf_{a\to\infty}q^{-1}\log[-\lambda_{\min}(T_i^E)]\ge c_*.
\tag{RBV25}
\]
All limits run over the original $a\equiv1\pmod4$ family at fixed
actual quartet/unit and eventually satisfy the explicitly stated
guards. No actual-period numerical eigenvalue is invented.

\subsection{Exact return through every original conductor chart}
This paragraph records the precise map needed to receive these
results on the supplied COB coefficient image, without identifying
two different sections of the same quotient. Write
$\chi(S)=\sum_{j=0}^qc_jS^j$, $c_q=1$, and define
\[
\mathfrak p_n(S)=\sum_{j=n+1}^qc_jS^{j-n-1},\qquad
\mathbb V_{n,\lambda}=\lambda^n,\qquad
\mathcal L\nu=\sum_\lambda\nu_\lambda\frac{\chi(S)}{S-\lambda},
\quad F=\mathcal L\mathbb V^{-1}.
\tag{RBV26}
\]
NCB5--7 proves $Fe_n=\mathfrak p_n$ by the coefficient of $S^{-1}$
in $\mathfrak p_nS^j/\chi$. Thus $F$ is invertible and its inverse is
the exact Laurent map
$(F^{-1}P)_n=[S^{-n-1}]P(S)/\chi(S)$, $0\le n<q$.
Let $C^{\mathsf T}$ be the original conductor multiplication in
exponential coefficients, and write
$\mathbb VC^{\mathsf T}=P_vJ$ with the original $v$ leading zeros.
For every retained invertible $q'$-row minor $J_I$, define
\[
\begin{gathered}
H_I=J_I^{-1}E_I^{\mathsf T},\quad
Q_I=E_{I^c}^{\mathsf T}-J_{I^c}H_I,\quad
L_I=FP_vE_{I^c},\quad C_E=\mathcal LC^{\mathsf T},\\
\eta_I=H_IP_v^{\mathsf T}F^{-1}P_a,\quad
z_I=Q_IP_v^{\mathsf T}F^{-1}P_a,\qquad
\boxed{P_a=C_E\eta_I+L_Iz_I.}
\end{gathered}\tag{RBV27}
\]
There is no low component because $\deg P_a<m$ makes all Laurent
coordinates below $q-m\ge v$ zero. To prove the displayed identity,
check the $I$ and $I^c$ rows to get
$JH_I+E_{I^c}Q_I=I$. Multiply by $FP_v$ and use
$FP_vJ=C_E$. These are the exact original residue and conductor maps;
no coefficient section is inserted in place of another.

In particular the consecutive minor
$I_0=\{0,\ldots,q'-1\}$ has determinant
\[
\det J_{I_0}=\mu_v^{q'}
 \prod_{r=0}^{q'-1}\binom{v+r}{v}\det\mathbb V_{\rm lower}\ne0.
\tag{RBV28}
\]
Indeed its row $r$ is the lower-root evaluation of a degree-$r$
polynomial with leading coefficient $\binom{v+r}{v}\mu_v$;
triangular coefficient elimination gives exactly this determinant.
Here $\eta_{I_0}=0$ and
$\operatorname{im}L_{I_0}=\mathcal P_{m-1}$ by the distinct monic
degrees $m-1,m-2,\ldots,0$ of its columns. The literal inverse is
\[
(z_{I_0})_j=[S^{-(q-m+j+1)}]P_a(S)/\chi(S).
\tag{RBV29}
\]
Therefore RBV8--25 apply to the Gram
$\Gamma_{i,I_0}=L_{I_0}^*G_i^EL_{I_0}$ and reflected matrix
$L_{I_0}^*W^EL_{I_0}$ with precisely the same bounds, because the
Rayleigh vector is the explicit column RBV29. For any actual other
chart $I$, RBV27 supplies the full original coefficient vector
$(\eta_I,z_I)$ instead. With $M_I=[C_E,L_I]$, the same calculation
holds on its full image with the two exact matrices
\[
M_I^*G_i^EM_I=\begin{pmatrix}
C_E^*G_i^EC_E&C_E^*G_i^EL_I\\
L_I^*G_i^EC_E&L_I^*G_i^EL_I
\end{pmatrix},\qquad M_I^*W^EM_I.
\tag{RBV30}
\]
The columns of $M_I$ form the original direct conductor/native
decomposition in the high-moment space, hence this Gram is positive
definite. All cross terms in RBV30 are essential.
This gives a proved negative class and exponential relative bound
through every complete original chart. Selecting only the native
block in a nonconsecutive chart is a further, different restriction:
RBV27 evaluates its omitted conductor component exactly instead of
declaring it zero. The companion original-image derivation identifies
which selector is the specified current COB map; no root ordering
alone is used here as a proof of an ordering of row subsets.
For the original four-copy maps use $I_4\otimes L_{I_0}$ or
$I_4\otimes M_I$ and the coefficient column with the displayed
one-copy preimage in its first block and zero in the other three.
Its reflected value and each boundary energy are exactly the
one-copy values. Hence RBV22--25 have the same constants on these
original four-copy images, with no additional rank factor.

\subsection{The complete tensor residue, its multiplier, and the residual}
We now expand the full FTR/FRT comparison on the present original
simple-quartet domain. For the four original roots
$z_{\epsilon\eta}=1/2+(2\epsilon-1)\delta+i(2\eta-1)\gamma$,
retain every original unit value in
\[
w_{\epsilon\eta}=\frac{v_h(z_{\epsilon\eta})}{h'(z_{\epsilon\eta})}\ne0,
\quad p(X,Y)=\sum_{\epsilon,\eta=0}^1w_{\epsilon\eta}X^\epsilon Y^\eta,
\quad c_{u,b}^{\rm ten}=[X^uY^b]p(X,Y)^a.
\tag{RBV31}
\]
There are no omitted jets: the one-factor multiplicity is exactly
one here. The original full-unit residue formula FTR.8, by taking
the simple residues in every ordered tensor factor, gives
\[
\mathcal B(f,g)=\sum_{u,b=0}^a
c_{u,b}^{\rm ten} f^{\sharp_a}(\lambda_{u,b})g(\lambda_{u,b}).
\tag{RBV32}
\]
Indeed each ordered root tuple contributes the product of its
$a$ unchanged $w_{\epsilon\eta}$ factors; the exponents of $X,Y$
record the two original plus counts. Expansion of $p^a$ counts
every ordered tuple once. Thus the tensor pullback multiplier in
the cyclic residue form is the actual class with values
$b_{\rm ten}(\lambda_{u,b})=\chi'(\lambda_{u,b})c_{u,b}^{\rm ten}$.
This is the complete simple-root specialization of FTR.6--12.

Put $\mathcal D=\{(u,b):c_{u,b}^{\rm ten}\ne0\}$ and
$\mathcal Z=\{0,\ldots,a\}^2\setminus\mathcal D$, and retain the
original idempotents
$e_{u,b}=\chi(S)/((S-\lambda_{u,b})\chi'(\lambda_{u,b}))$.
The trace multiplier and complementary form are exactly
\[
\begin{gathered}
a_{\rm tr}=\sum_{(u,b)\in\mathcal D}
              (c_{u,b}^{\rm ten})^{-1}e_{u,b},\qquad
\mathcal W_{\mathcal Z}(f,g)=\sum_{(u,b)\in\mathcal Z}
              f^{\sharp_a}(\lambda_{u,b})g(\lambda_{u,b}),\\
b_{\rm ten}a_{\rm tr}=e_{\mathcal D}\chi'\quad\text{in }E,
\qquad
\boxed{\mathcal W(f,g)=\mathcal B(f,a_{\rm tr}g)
                              +\mathcal W_{\mathcal Z}(f,g).}
\end{gathered}\tag{RBV33}
\]
All equations follow by evaluation at every simple root; the
evaluation map is an isomorphism, so they prove the typed algebra
map, not just matching dimensions. They are exactly FRT.2--5.
The trace equality can also be verified by taking the residue of
$\chi'f^{\sharp_a}g/\chi$ at each root. The classical multiplication
trace construction is recorded by the Stacks Project Authors,
Section 49.3, tags 0BVH and 0BJF \cite{human:stacks-trace}; the
displayed root computations prove the full specialization here.

In particular the retained residual has the exact finite value
\[
\mathcal W_{\mathcal Z}(P_a,P_a)
 =\sum_{(u,b)\in\mathcal Z}P_a(\lambda_{u,b})^2.
\tag{RBV34}
\]
Conjugation of the full unit and $h$ gives
$w_{\epsilon,1-\eta}=\overline{w_{\epsilon\eta}}$;
their original reflection gives
$w_{1-\epsilon,\eta}=-\overline{w_{\epsilon\eta}}$.
Since $a$ is odd, the coefficient sets $\mathcal D,\mathcal Z$
are closed under both grid symmetries. Thus both partial sums of
$P_a(\lambda)^2$ are real. RBV6 proves that RBV34 is nonpositive,
and strictly negative whenever $\mathcal Z$ is nonempty. This
does not determine that set by a rank assumption: its exact definition
retains all cancellations in the actual coefficients RBV31.

There is also a negative pairing directly in the retained tensor
residue without needing to determine its interior zero coefficients.
The four edge coefficients of $p^a$ are binomial products of
nonzero original $w$'s, so every perimeter point belongs to
$\mathcal D$. In particular its two complete edges $u=0,a$
contain $2(a+1)$ points with absolute comparison value
$\sinh(r_0a)$. Applying RBV8a at these points and RBV6 at all
remaining points of $\mathcal D$ proves
\[
\boxed{\mathcal B(P_a,a_{\rm tr}P_a)
 \le-2(a+1)(\sinh(r_0a)-\epsilon_a)^2<0.}\tag{RBV35}
\]
The second argument is the explicitly computed full polynomial
class $a_{\rm tr}P_a\bmod\chi$; it is not identified with a native
coefficient vector unless its Laurent coordinates satisfy that
specific section's membership test.

Finally the original arithmetic injection is
$\mathcal I:E\hookrightarrow Q^{\otimes a}$ with the complete
$[v_h]_h^{\otimes a}$ factor. Its inverse on its named image is
the original full Mellin-jet observation, division by that invertible
Taylor factor and inverse sum substitution. Thus
$\mathcal I M_{a_{\rm tr}}\mathcal I^{-1}$ is the exact arithmetic
map realizing RBV33--35 on $\mathcal I(e_{\mathcal D}E)$.
In RBV35 the first argument may be written $e_{\mathcal D}P_a$,
since the complementary first-input block pairs to zero under
$\mathcal B$ by RBV32; the separate value RBV34 remains in
the full reflected trace.
It commutes with the original sum action because both act by
multiplication before this injection. The proper boundary image
uses instead the full COB.8 map
$\Psi_2:[P]\mapsto[P(D_1+\cdots+D_a)F_h^{\otimes a}]$
on $\mathbb C[S]/\chi^{a+1}$, whose one-factor Taylor unit is
$[v_h]_{h^2}$; these two different original target spaces are not
silently identified. The ordinary class is recovered on the named
boundary image by $\mathcal I\Delta_i^{-1}\Psi_2^{-1}$ exactly
as in COB.12. This retains the full tensor-residue residual and
every original unit. The finite packet trace is not substituted
for the full infinite Weil distribution, and none of these values
is assigned to the native minimum metric.

The tensor calculation also has a relative-metric receiver, with
the residual retained at matrix level. On any of the certified
images above let $L$ denote its displayed injection and set
$\Gamma_i=L^*G_i^EL$. With $W_{\mathcal D}^E,W_{\mathcal Z}^E$
the two explicitly computed reflected matrices, define
\[
\begin{gathered}
T_{i,\mathcal D}=\Gamma_i^{-1/2}L^*W_{\mathcal D}^EL\Gamma_i^{-1/2},
\qquad T_{i,\mathcal Z}=\Gamma_i^{-1/2}L^*W_{\mathcal Z}^EL\Gamma_i^{-1/2},\\
T_i=T_{i,\mathcal D}+T_{i,\mathcal Z},\qquad
\boxed{\lambda_{\min}(T_{i,\mathcal D})
\le-\frac{2(a+1)(\sinh(r_0a)-\epsilon_a)^2}{B_a}<0.}
\end{gathered}\tag{RBV36}
\]
Both partial matrices are Hermitian by their proved symmetry sets.
Use the very same preimage of $P_a$ as in RBV22 and the strict
boundary bound RBV21 to prove the displayed inequality. Its
negative magnitude again has logarithmic lower rate $c_*$ after
division by $q$. This is a bound for the actual tensor-residue
comparison $\mathcal B(f,a_{\rm tr}g)$; the other matrix remains
in the exact displayed sum. No equality of the smallest eigenvalues
of those three matrices is asserted.
\endgroup
